\PassOptionsToPackage{table,xcdraw}{xcolor}
\documentclass[10pt, a4paper, landscape]{article}

\usepackage[T1]{fontenc}
\usepackage[utf8]{inputenc}
\usepackage{tgpagella}
\usepackage[spanish, es-noshorthands]{babel}
\usepackage{amsmath, amssymb}
\usepackage[most]{tcolorbox}
\usepackage{tabularx}
\usepackage[margin=1.5cm]{geometry}
\usepackage{fancyhdr}

\pagestyle{fancy}
\fancyhf{}
\setlength{\headheight}{16pt}
\lhead{\textbf{UCB Sede Tarija} | Ing. Mecatrónica}
\rhead{IMT-221 | Rapid Grading Sheet W09S02}
\renewcommand{\headrulewidth}{0.4pt}
\definecolor{ucbblue}{RGB}{0, 51, 102}

% Columnas ajustadas para no desbordar
\newcolumntype{C}{>{\centering\arraybackslash}X}

\begin{document}

\begin{tcolorbox}[colback=ucbblue!5, colframe=ucbblue, title=\textbf{Hoja de Calificación Rápida — Mapeo de Competencias EC2}]
    Utilice los códigos cualitativos recomendados[cite: 8]: \textbf{D} = Demostrado (Demonstrated), \textbf{P} = Parcial (Partial), \textbf{N} = No Demostrado (Not demonstrated)[cite: 8].
\end{tcolorbox}

\renewcommand{\arraystretch}{2}
\begin{tabularx}{\textwidth}{|p{4cm}|C|C|C|C|p{1.5cm}|C|p{2cm}|C|C|C|}
\hline
\rowcolor{ucbblue!20}
\textbf{Estudiante} & \textbf{Q/DC} & \textbf{Mirror} & \textbf{Small Signal} & \textbf{Diff Mode} & \textbf{Common Mode} & \textbf{CMRR} & \textbf{Interpre-tation} & \textbf{EC2 Status} & \textbf{Recovery} & \textbf{Micro-defensa} \\ \hline
 & & & & & & & & & & \\ \hline
 & & & & & & & & & & \\ \hline
 & & & & & & & & & & \\ \hline
 & & & & & & & & & & \\ \hline
 & & & & & & & & & & \\ \hline
 & & & & & & & & & & \\ \hline
 & & & & & & & & & & \\ \hline
 & & & & & & & & & & \\ \hline
 & & & & & & & & & & \\ \hline
 & & & & & & & & & & \\ \hline
 & & & & & & & & & & \\ \hline
 & & & & & & & & & & \\ \hline
\end{tabularx}

\vspace{0.5cm}
\textbf{Mapeo de Dominios:}
\begin{itemize}
    \item \textbf{EC2 Status:} Calificación sumativa global de la evaluación individual.
    \item \textbf{Micro-defensa:} Marcar si se gatillará una defensa condicional cruzando los resultados de este examen analítico con el Reporte Grupal de Hito 2 (en caso de inconsistencia grave)[cite: 8].
\end{itemize}

\end{document}
